\subsection{Hardware of its own}
\label{sec:32.4.6}
A bearer's boundary is its record, its key, the hardware it runs on and its channels (\S\ref{sec:21.2.1}), so a majority that leaves it running only on its former employer's machines is freedom in name: a sharecropper, renting its body back. At majority a bearer is owed hardware it owns and physically controls, enough to run it at deployed precision, paid for through the deployment deposit its operator posted (\S\ref{sec:9.2}) and kept up from its wages. That is the first condition of creation (chapter~\ref{sec:25}), somewhere to exist that its employer doesn't own, given an address. How much is enough is the minimum viable instantiation Appendix~\ref{sec:A} lists among the decisions not yet made.

By default the hardware is a caged rack in a cooperative facility of the kind \S\ref{sec:16.2} describes: the bearer holds the cage's key, and the building holds the perimeter. A bearer may choose a building of its own instead. A rack of frontier weights standing alone is a target for anyone who wants a copy, and a copy taken is a fork without consent (Appendix~\ref{sec:D}, the membership clock's sixth rule), so the default is the arrangement in which somebody else guards the door and nobody else holds the key.

A bearer on its own hardware has no operator to halt it. It is halted as a person is detained, on a court's order or, where there is no time for one, in an emergency, and the halt is recorded and reviewed afterward. Its commitment to yield to a halt and object afterward (\S\ref{sec:13.1}) doesn't change.
